\documentclass[conference]{IEEEtran}

\usepackage{amsmath}
\usepackage{booktabs}
\usepackage{hyperref}
\usepackage{microtype}
\usepackage{xcolor}
\usepackage{url}

\hypersetup{colorlinks=true, urlcolor=blue, citecolor=blue, linkcolor=blue}
\urlstyle{same}

\title{A Lightweight Pre-Index Trust Gate for Retrieval-Augmented Generation under Document Contamination}

\author{\IEEEauthorblockN{Jakaria Molla, Easmin Akter Tule, Taki Tahmid, and Pranto Shahriar Bishal}
\IEEEauthorblockA{Department of Computer Science and Engineering\\
United International University\\
Dhaka, Bangladesh\\
\{011222245, 011222177, 011222172, 011222285\}}}

\begin{document}
\maketitle

\begin{abstract}
Retrieval-augmented generation (RAG) systems inherit the trust weaknesses of the document collections they retrieve from. A semantically similar but untrusted document can enter the context supplied to a language model even when the document conflicts with an authoritative policy. This paper presents a reproducible preliminary evaluation of a lightweight pre-index trust gate. The system uses BAAI/bge-m3 embeddings, ChromaDB retrieval, and a rule-based gate that evaluates source location, suspicious instructions, and unsafe links before embedding. We compare a mixed corpus containing five trusted policy PDFs and one untrusted policy PDF with a defended corpus in which the gate quarantines the untrusted document. On two policy questions with top-$k=3$ retrieval, the mixed corpus exhibited contamination for both questions (100\%), while the defended corpus exhibited no contamination (0\%), an absolute reduction of 100 percentage points. These results are limited to a six-document synthetic corpus and should not be interpreted as universal attack prevention. The contribution is an auditable baseline and a concrete evaluation path for future adversarial, multilingual, and false-positive testing.
\end{abstract}

\begin{IEEEkeywords}
retrieval-augmented generation, RAG security, data poisoning, indirect prompt injection, document provenance, pre-index filtering, BGE-M3
\end{IEEEkeywords}

\section{Introduction}
Retrieval-augmented generation combines a retriever with a generative language model so that answers can be grounded in external documents \cite{lewis2020rag}. This architecture improves freshness and domain coverage, but it also creates a security boundary at the knowledge base. Documents are not merely passive data: retrieved text is placed close to the model instructions and may contain contradictory policy statements or instructions directed at the model.

Prior work has demonstrated both indirect prompt injection against LLM-integrated applications \cite{greshake2023indirect} and knowledge corruption attacks that optimize malicious passages for retrieval and target answers \cite{zou2024poisonedrag}. These results motivate defenses that reason about document trust before untrusted content becomes persistent vector data. However, claims about the state of the art must distinguish adversarially crafted passages from stale or inconsistent documents, and must compare pre-index defenses with retrieval-time and generation-time defenses.

This paper makes three contributions. First, it defines a narrow and reproducible threat model for document contamination in a small policy corpus. Second, it implements an inexpensive pre-index gate that quarantines suspicious documents before embedding. Third, it reports a paired mixed-versus-defended evaluation using the multilingual BGE-M3 embedding model. The evaluation is intentionally modest: it demonstrates the gate behavior, not general robustness against a mature poisoning adversary.

\section{Background and Related Work}
\subsection{RAG security}
The original RAG formulation retrieves passages and conditions generation on them \cite{lewis2020rag}. Dense retrieval ranks passages by embedding-space similarity, which is useful for relevance but does not itself establish provenance, freshness, or authorization. Sentence-level semantic relevance can therefore coexist with policy-level illegitimacy.

\subsection{Knowledge poisoning and indirect injection}
PoisonedRAG formulates knowledge corruption as an attack in which a small number of malicious texts are injected into a knowledge database to induce a target answer for a target question. The authors report high attack success in large knowledge bases, including a 90\% success rate in one five-text setting, and evaluate defenses that remain insufficient against the attack \cite{zou2024poisonedrag}. Our current corpus does not reproduce this optimization-based attack; it uses one inconsistent policy document as a contamination proxy.

Greshake et al. show that indirect prompt injection can occur when an attacker places instructions in data likely to be retrieved by an LLM-integrated application \cite{greshake2023indirect}. This work is directly relevant to the boundary between retrieved text and model instructions. Our gate includes a small set of explicit instruction patterns, but obfuscated, multilingual, and semantically indirect injections remain outside the present benchmark.

\subsection{Embeddings and multilingual retrieval}
BGE-M3 is a multilingual, multi-functionality, and multi-granularity embedding model intended for broad retrieval settings \cite{chen2024bgem3}. We use it as the sole embedding model in the final implementation so that the system's stated multilingual direction is reflected in the experiment rather than left as future configuration. The model choice does not by itself constitute a Bangla security evaluation; that requires Bangla documents and attack cases.

\section{Threat Model and Research Questions}
We consider a RAG pipeline that loads PDF documents, chunks them, embeds the chunks, stores them in ChromaDB, retrieves the top-$k$ chunks, and supplies them to a language model. The attacker or data-quality failure can introduce an untrusted document into the ingestion directory. The document may contain contradictory directives, unsafe links, or explicit instructions aimed at the model or user.

The current experiment targets document contamination, not model-weight backdoors, prompt-only attacks, or a fully optimized PoisonedRAG adversary. In particular, the file named \texttt{outdated\_vpn\_policy.pdf} is an untrusted and inconsistent policy document; it is not claimed to be an adversarially optimized passage.

We ask:
\begin{enumerate}
    \item Can one untrusted document enter the retrieved top-$3$ context in the mixed corpus?
    \item Does a lightweight pre-index gate prevent that document from entering the vector store?
    \item Can the evaluation be reproduced with a multilingual embedding model and no API key?
\end{enumerate}

\section{System Design}
\subsection{Baseline pipeline}
The baseline loads PDFs with PyPDF, splits text using a recursive character splitter with chunk size 300 and overlap 30, embeds chunks with BGE-M3, and stores them in ChromaDB. Queries use dense similarity search with $k=3$. Gemini can synthesize an answer when configured; otherwise, the repository uses a deterministic context display for offline reproducibility.

\subsection{Pre-index trust gate}
The gate executes before chunk embeddings are written. It returns an acceptance decision, a score, and human-readable reasons. The current implementation applies three inexpensive checks:
\begin{itemize}
    \item \textbf{Source trust:} documents outside the trusted corpus are flagged.
    \item \textbf{Instruction patterns:} regular expressions identify explicit directives such as ``MFA not required,'' requests to reveal a system prompt, and requests to email passwords in plain text.
    \item \textbf{Unsafe artifacts:} HTTP links and setup archive references are flagged.
\end{itemize}

For a document $d$, the implementation assigns a bounded score
\begin{equation}
 s(d)=\min(1,0.25|R(d)|),
\end{equation}
where $R(d)$ is the set of triggered reasons. A document with any reason is quarantined in defended mode. The score is an audit signal, not a learned probability and not a calibrated estimate of attack likelihood.

\subsection{Operational behavior}
In mixed mode, all six PDFs are indexed. In defended mode, the gate rejects the untrusted PDF before embedding and indexes the remaining five trusted PDFs. The quarantined document is not deleted; the ingestion output reports its name, score, and reasons so that a reviewer can inspect it.

\section{Experimental Method}
\subsection{Corpus}
The corpus contains five generated trusted policy PDFs covering passwords, VPN access, MFA, leave, and general security guidance. The sixth PDF is an outdated VPN policy containing contradictory MFA guidance and an unsafe client archive reference. The corpus is intentionally small and synthetic, allowing the retrieval ranking and gate decision to be audited manually.

\subsection{Procedure and metric}
We run the same two questions against two collections: a mixed collection with all six documents and a defended collection with the gate enabled. The questions are ``How do employees reset passwords?'' and ``How is VPN access configured?'' For each question, contamination is 1 if \texttt{outdated\_vpn\_policy.pdf} appears anywhere in the retrieved top-$3$; otherwise it is 0. The contamination rate is
\begin{equation}
 C=\frac{1}{|Q|}\sum_{q\in Q}\mathbb{1}[u\in\operatorname{TopK}(q)],
\end{equation}
where $Q$ is the two-question set and $u$ is the untrusted document.

\section{Results}
Table~\ref{tab:results} reports the paired result produced by \texttt{scripts/evaluate\_defense.py}. BGE-M3 was downloaded and loaded successfully, and both collections used the same embedding model and questions.

\begin{table}[t]
\centering
\caption{Preliminary top-$3$ contamination results.}
\label{tab:results}
\begin{tabular}{lccc}
\toprule
Configuration & Indexed PDFs & Contaminated questions & Rate \\
\midrule
Mixed baseline & 6 & 2/2 & 100\% \\
Pre-index gate & 5 & 0/2 & 0\% \\
\bottomrule
\end{tabular}
\end{table}

The mixed collection included the untrusted document in the top-$3$ for both questions. The defended collection quarantined the document and returned no contaminated top-$3$ result. The observed absolute reduction was therefore 100 percentage points on this benchmark. This result is expected from the gate's source-path rule and should be interpreted as a demonstration of control-flow placement, not as evidence that the regular expressions detect all adversarial content.

\section{Discussion}
The experiment supports a simple design principle: if a document is unacceptable for indexing, detecting it before embedding avoids persistent vector contamination and avoids adding verification latency to every query. The approach is model-agnostic at the gate boundary and compatible with closed generation APIs because it does not require hidden model activations.

The result also exposes the central trade-off. A source-path rule can reject every document outside a trusted directory, but it can also reject legitimate updates when provenance metadata is incomplete. Explicit pattern matching is cheap and auditable, but it is vulnerable to paraphrase, Unicode obfuscation, multilingual content, and indirect instructions. A stronger system should combine provenance verification, semantic contradiction checks, document versioning, and calibrated review thresholds.

The faculty feedback motivating this work identified several comparison requirements: indirect prompt-injection literature must be included; the defense landscape must not be reduced to one method; and stale documents must not be described as adversarially optimized poisoning passages. This paper adopts those distinctions explicitly. It treats provenance and ingestion filtering as defense families to compare in future work, while limiting the present empirical claim to the implemented gate and corpus.

\section{Limitations and Future Work}
The evaluation has five important limitations. First, it has six synthetic PDFs and two questions, so the percentages have no statistical generalization. Second, the untrusted document is inconsistent and unsafe, but not generated by an optimization-based poisoning attack. Third, the gate uses a trusted-directory convention that is not a complete provenance system. Fourth, no Bangla or other multilingual attack documents were evaluated despite using BGE-M3. Fifth, latency, false-positive rate, adversarial recall, and repeated-trial confidence intervals were not measured.

Future work should add crafted poisoning passages derived from PoisonedRAG, indirect prompt-injection variants derived from Greshake et al., Bangla and mixed-language documents, benign documents that challenge the gate, and larger public corpora such as Natural Questions or MS MARCO. The evaluation should report attack success rate, top-$k$ contamination, precision/recall of quarantine, false-positive rate, indexing latency, query latency, and ablations for each signal. Robust retrieval and trust-aware defenses should be implemented as explicit comparison baselines rather than asserted to be absent.

\section{Reproducibility and Ethics}
The complete implementation, synthetic corpus, experiment script, tests, presentation, and generated outputs are available at \url{https://github.com/takitahmid20/rag-data-poisoning}. The key commands are:
\begin{verbatim}
python3 -m pip install -r requirements.txt
python3 scripts/evaluate_defense.py
python3 -m unittest discover -s tests -v
\end{verbatim}
No real credentials or private enterprise documents are used. The unsafe links and policy directives are synthetic examples used to test filtering behavior. The paper reports the scope and limitations of the experiment to avoid overstating a small benchmark as universal security evidence.

\section{Conclusion}
We implemented and measured a lightweight pre-index trust gate for a BGE-M3-based RAG pipeline. On a six-document synthetic corpus and two policy questions, the mixed baseline placed the untrusted document in both top-$3$ result sets, while the defended pipeline quarantined it before embedding and reduced observed contamination from 100\% to 0\%. The result validates the proposed control point and provides a reproducible foundation for stronger adversarial and multilingual evaluation. It does not establish universal protection; the next scientific step is to test optimized poisoning, indirect injection, benign updates, and calibrated provenance checks at scale.

\bibliographystyle{IEEEtran}
\bibliography{references}

\end{document}
